\chapter*{Abstract}
\addcontentsline{toc}{chapter}{Abstract}

Modern robotic and edge computing systems increasingly execute computer-vision workloads on
embedded platforms, where compute, memory, and energy budgets are
constrained. 

The ROS~2 middleware promotes modular architectures through
nodes that communicate by message passing; however, each communication
boundary between nodes can introduce serialization, memory-transfer, and
scheduling overheads that affect latency and throughput, an effect that is
particularly relevant on resource-limited platforms. 

This work experimentally characterizes the performance impact of partitioning a video
processing pipeline into ROS~2 nodes (rclpy), compared with an equivalent
single-process Python implementation, on an NVIDIA Jetson Orin Nano.

Functionally equivalent implementations of the same computer-vision
application, based on the YOLOv8n object detector, are evaluated under a
reproducible
benchmarking methodology, with increasing partitioning configurations and several input resolutions,
measuring end-to-end latency, throughput (FPS), and CPU, RAM, and GPU
utilization.

\todo{RESULTS: summarize the main finding once the experiments are complete
(latency, throughput, and resource-utilization differences between
configurations, and the effect of the degree of partitioning).}

\todo{CONCLUSION: close with the answer to the research question and the
practical implication for developers of embedded robotic systems.} The
primary contribution is a controlled, quantitative characterization of the
trade-off between the modularity offered by ROS~2 and its performance cost
for this class of workload on embedded platforms.

\bigskip
\noindent\textbf{Keywords:} ROS~2, computer vision, embedded platforms,
NVIDIA Jetson Orin Nano, performance, latency.
